\chapter{Chuỗi Markov thời gian rời rạc\\{\normalsize\textit{Discrete-Time Markov Chains}}}
\label{ch:discrete}

% ============================================================================
\section{Tại sao chuỗi Markov quan trọng? (Why Markov Chains Matter)}
\label{sec:why-markov}

\textbf{Câu chuyện:} Hãy tưởng tượng bạn đang chơi cờ tỷ phú (Monopoly).
Bạn đang ở ô ``Nhà ga''. Bạn sẽ đi đâu tiếp theo?
Điều đó chỉ phụ thuộc vào ô hiện tại và kết quả gieo xúc xắc ---
không quan trọng bạn đã đi qua những ô nào trước đó.
Đây chính là ý tưởng cốt lõi của chuỗi Markov: \textbf{tương lai chỉ phụ thuộc hiện tại, không phụ thuộc quá khứ}.

\textbf{Ứng dụng thực tế:}
\begin{itemize}[nosep]
  \item \textbf{Thời tiết:} Ngày mai mưa hay nắng chỉ phụ thuộc hôm nay, không phụ thuộc tuần trước.
  \item \textbf{Xếp hạng Google (PageRank):} Người lướt web ``nhảy'' ngẫu nhiên giữa các trang.
  \item \textbf{Tài chính:} Giá cổ phiếu ngày mai phụ thuộc giá hôm nay (mô hình đơn giản).
  \item \textbf{Sinh học:} Kích thước quần thể thế hệ sau phụ thuộc thế hệ hiện tại.
  \item \textbf{Truyền thông:} Xếp hạng tín dụng năm sau phụ thuộc xếp hạng năm nay.
\end{itemize}

% ============================================================================
\section{Tính chất Markov (Markov Property)}
\label{sec:markov-property}

\textbf{Câu chuyện:} Hãy tưởng tượng một con kiến trên bàn cờ. Con kiến chỉ biết ô nó đang đứng.
Nó không nhớ đã đi qua bao nhiêu ô, cũng không nhớ con đường đã đi.
Quyết định bước tiếp theo chỉ dựa vào vị trí hiện tại --- đây là ``con kiến Markov''.

\begin{dinhnghia}[(Tính chất Markov -- \en{Markov property})]
Một \term{quá trình ngẫu nhiên} (\en{stochastic process}) $(Z_n)_{n\in\N}$ nhận giá trị trong $\Z$
được gọi là \term{Markov}, hay có \term{tính chất Markov}, nếu với mọi $n \geq 1$,
phân phối xác suất của $Z_{n+1}$ chỉ được xác định bởi trạng thái $Z_n$ tại thời điểm $n$,
và không phụ thuộc vào các giá trị quá khứ $Z_k$ với $k \leq n-1$.

Nói cách khác, với mọi $n \geq 1$ và mọi $j, i_n, \ldots, i_1 \in \Z$:
\[
  \PP(Z_{n+1} = j \mid Z_n = i_n,\; Z_{n-1} = i_{n-1},\; \ldots,\; Z_0 = i_0)
  = \PP(Z_{n+1} = j \mid Z_n = i_n).
\]
\end{dinhnghia}

\begin{tomtat}
Nói đơn giản: nếu bạn biết trạng thái hiện tại $Z_n$,
thì biết thêm quá khứ $(Z_0, Z_1, \ldots, Z_{n-1})$ cũng không giúp dự đoán tương lai $Z_{n+1}$ tốt hơn.
Đây chính là tính chất ``không nhớ'' (\en{memoryless property}).

Ví dụ con kiến: biết con kiến đang ở ô $A$ là đủ. Biết thêm ``nó đi từ $C \to B \to A$'' không thay đổi xác suất bước tiếp theo.
\end{tomtat}

Đặc biệt ta có:
\[
  \PP(Z_{n+1} = j \mid Z_n = i_n,\; Z_{n-1} = i_{n-1}) = \PP(Z_{n+1} = j \mid Z_n = i_n),
\]
và:
\[
  \PP(Z_2 = j \mid Z_1 = i_1,\; Z_0 = i_0) = \PP(Z_2 = j \mid Z_1 = i_1).
\]

Mặt khác, các \term{xác suất chuyển bậc một} (\en{first order transition probabilities})
có thể dùng để tính toàn bộ phân phối của quá trình bằng quy nạp:
\begin{align}
  \PP(Z_n = i_n, Z_{n-1} = i_{n-1}, \ldots, Z_0 = i_0)
  &= P_{i_{n-1}, i_n} \cdots P_{i_0, i_1} \, \PP(Z_0 = i_0),
\end{align}
với $i_0, i_1, \ldots, i_n \in \Z$.

Theo \term{quy tắc xác suất toàn phần} (\en{law of total probability}):
\begin{equation}\label{eq:total-prob}
  \PP(Z_1 = i) = \sum_{j \in \Z} \PP(Z_1 = i \mid Z_0 = j)\,\PP(Z_0 = j), \qquad i \in \Z.
\end{equation}

% --- Example: Random walk ---
\begin{example}[Bước đi ngẫu nhiên -- \en{Random walk}]
Bước đi ngẫu nhiên (\en{random walk}):
\begin{equation}\label{eq:random-walk}
  S_n = X_1 + \cdots + X_n, \qquad n \in \N,
\end{equation}
trong đó $(X_n)_{n \geq 1}$ là dãy các biến ngẫu nhiên \term{độc lập} (\en{independent})
nhận giá trị trong $\Z$, là một chuỗi Markov thời gian rời rạc.

\textbf{Chứng minh:} Ta cần chỉ ra tính Markov. Với mọi $j, i_n, \ldots, i_1 \in \Z$:
\begin{align*}
  \PP(S_{n+1} = j \mid S_n = i_n,\; S_{n-1} = i_{n-1},\; \ldots,\; S_1 = i_1)
  &= \PP(S_n + X_{n+1} = j \mid S_n = i_n, \ldots) \\
  &= \PP(X_{n+1} = j - i_n \mid S_n = i_n, \ldots) \\
  &= \PP(X_{n+1} = j - i_n) \\
  &= \PP(S_{n+1} = j \mid S_n = i_n).
\end{align*}
Bước thứ ba dùng tính \term{độc lập} của $X_{n+1}$ với $S_1, \ldots, S_n$ (vì $X_{n+1}$ độc lập với $X_1, \ldots, X_n$).
\end{example}

Tổng quát hơn, mọi quá trình có \term{gia số độc lập} (\en{independent increments}) đều là chuỗi Markov. Tuy nhiên, không phải mọi chuỗi Markov đều có gia số độc lập.

% ============================================================================
\subsection{Ma trận chuyển (Transition Matrix)}
\label{subsec:transition-matrix}

\textbf{Câu chuyện:} Ma trận chuyển là ``bản đồ xác suất'': từ mỗi thành phố (trạng thái),
bạn biết xác suất đi đến các thành phố khác. Hàng $i$ trong ma trận cho bạn biết:
``Nếu đang ở thành phố $i$, xác suất đi đến mỗi thành phố khác là bao nhiêu?''

Sự tiến hóa ngẫu nhiên của quá trình Markov $(Z_n)_{n\in\N}$ được xác định bởi dữ liệu:
\[
  P_{i,j} := \PP(Z_{n+1} = j \mid Z_n = i), \qquad i, j \in \Z,
\]
trong đó ta giả sử xác suất $\PP(Z_{n+1} = j \mid Z_n = i)$ không phụ thuộc vào $n \in \N$.
Khi đó chuỗi Markov $(Z_n)_{n\in\N}$ được gọi là \term{đồng nhất theo thời gian} (\en{time homogeneous}).

\begin{dinhnghia}[(Ma trận chuyển -- \en{Transition matrix})]
Dữ liệu trên được mã hóa thành một ma trận, gọi là \term{ma trận chuyển} (\en{transition matrix}) của chuỗi Markov:
\[
  [P_{i,j}]_{i,j\in\Z} = [\PP(Z_{n+1} = j \mid Z_n = i)]_{i,j\in\Z}.
\]
\end{dinhnghia}

\textbf{Lưu ý quan trọng:} Thứ tự các chỉ số $(i,j)$ trong $\PP(Z_{n+1} = j \mid Z_n = i)$ và $P_{i,j}$ bị đảo ngược. Cụ thể, trạng thái ban đầu $i$ là \emph{số hàng} trong ma trận, còn trạng thái đích $j$ là \emph{số cột}.

\begin{dinhnghia}[(Trạng thái hấp thụ -- \en{Absorbing state})]
Trạng thái $k \in \Z$ được gọi là \term{hấp thụ} (\en{absorbing}) nếu $P_{k,k} = 1$.
\end{dinhnghia}

Do quan hệ:
\begin{equation}\label{eq:row-sum-1}
  \sum_{j \in \Z} \PP(Z_{n+1} = j \mid Z_n = i) = 1, \qquad i \in \N,
\end{equation}
các \emph{hàng} của ma trận chuyển thỏa mãn: $\sum_{j} P_{i,j} = 1$ với mọi $i$.

% --- Finite state space ---
Trong trường hợp chuỗi Markov $(Z_k)_{k\in\N}$ nhận giá trị trong không gian trạng thái hữu hạn $\{0, \ldots, N\}$, ma trận chuyển có dạng:
\[
  [P_{i,j}]_{0 \leq i,j \leq N}
  = \begin{bmatrix}
    P_{0,0} & P_{0,1} & P_{0,2} & \cdots & P_{0,N} \\
    P_{1,0} & P_{1,1} & P_{1,2} & \cdots & P_{1,N} \\
    P_{2,0} & P_{2,1} & P_{2,2} & \cdots & P_{2,N} \\
    \vdots  & \vdots  & \vdots  & \ddots & \vdots \\
    P_{N,0} & P_{N,1} & P_{N,2} & \cdots & P_{N,N}
  \end{bmatrix}.
\]

% --- Examples ---
\begin{example}[Quá trình cờ bạc -- \en{Gambling process}]
\label{ex:gambler}
Một người chơi bắt đầu với $i$ đô la. Mỗi ván, anh ta thắng $1$ đô la (xác suất $p$)
hoặc thua $1$ đô la (xác suất $q = 1-p$). Trò chơi dừng khi cháy túi ($0$ đô la) hoặc đạt mục tiêu ($N$ đô la).

Ma trận chuyển trên $\{0, 1, \ldots, N\}$ với các trạng thái hấp thụ $0$ và $N$:
\[
  P = [P_{i,j}]_{0 \leq i,j \leq N}
  = \begin{bmatrix}
    1   & 0 & 0 & \cdots & 0 & 0 & 0 \\
    q   & 0 & p & \cdots & 0 & 0 & 0 \\
    \vdots & \vdots & \vdots & \ddots & \vdots & \vdots & \vdots \\
    0   & 0 & 0 & \cdots & q & 0 & p \\
    0   & 0 & 0 & \cdots & 0 & 0 & 1
  \end{bmatrix}.
\]
\end{example}

\begin{example}[Chuỗi Ehrenfest -- \en{Ehrenfest chain}]
\label{ex:ehrenfest}
Hai bình khí (trái và phải) nối với nhau bằng một lỗ, chứa tổng cộng $N$ viên bi.
Tại mỗi bước, ta chọn ngẫu nhiên một viên bi và chuyển nó sang bình kia.
Gọi $X_n \in \{0,\ldots,N\}$ là số bi bên trái tại thời điểm $n$.
Xác suất chuyển:
\[
  \PP(X_{n+1} = k+1 \mid X_n = k) = \frac{N-k}{N}, \qquad
  \PP(X_{n+1} = k-1 \mid X_n = k) = \frac{k}{N}.
\]

\textbf{Ý nghĩa vật lý:} Mô hình này mô phỏng sự khuếch tán khí giữa hai ngăn.
Nếu bên trái có nhiều bi hơn ($k > N/2$), xác suất giảm ($k/N$) lớn hơn xác suất tăng ($(N-k)/N$)
--- hệ thống ``kéo'' về trạng thái cân bằng $k = N/2$.
\end{example}

\begin{example}[Ehrenfest $N = 4$ -- ma trận cụ thể]
\label{ex:ehrenfest-N4}
Với $N = 4$ viên bi, ma trận chuyển $5 \times 5$:
\[
  P = \begin{bmatrix}
    0 & 1 & 0 & 0 & 0 \\
    1/4 & 0 & 3/4 & 0 & 0 \\
    0 & 2/4 & 0 & 2/4 & 0 \\
    0 & 0 & 3/4 & 0 & 1/4 \\
    0 & 0 & 0 & 1 & 0
  \end{bmatrix}.
\]

\textbf{Đọc ma trận:}
\begin{itemize}[nosep]
  \item Hàng 0: nếu bên trái có 0 bi, chắc chắn tăng lên 1 (chọn bi nào cũng từ bên phải).
  \item Hàng 2: nếu bên trái có 2 bi, xác suất giảm = $2/4$, xác suất tăng = $2/4$ (cân bằng!).
  \item Hàng 4: nếu bên trái có 4 bi (tất cả), chắc chắn giảm xuống 3.
\end{itemize}
\end{example}

\begin{example}[Xếp hạng tín dụng -- \en{Credit rating}]
Trong tài chính, xếp hạng tín dụng (AAA, AA, A, BBB, \ldots, D) có thể được mô hình hóa
bằng chuỗi Markov. Ma trận chuyển ghi lại xác suất thay đổi hạng tín dụng sau mỗi năm.
Các hạng cao hơn có hệ số đường chéo lớn hơn (ổn định hơn).
\end{example}

% ============================================================================
\subsection{Xác suất chuyển nhiều bước (Multi-step Transition Probabilities)}

Ta muốn tính \term{xác suất chuyển hai bước} (\en{two-step transition probability})
$\PP(Z_{n+2} = j \mid Z_n = i)$, không có trực tiếp trong ma trận chuyển.

Bằng \term{phân tích bước đầu tiên} (\en{first step analysis}), ký hiệu $S$ là không gian trạng thái:
\begin{align*}
  \PP(Z_{n+2} = j \mid Z_n = i)
  &= \sum_{l \in S} \PP(Z_{n+1} = l \mid Z_n = i)\,\PP(Z_{n+2} = j \mid Z_{n+1} = l)\\
  &= \sum_{l \in S} P_{i,l}\,P_{l,j}
  = P^2_{i,j}, \qquad i, j \in S.
\end{align*}

Tổng quát hơn, với mọi $k \in \N$:

\begin{congthuc}
\textbf{Xác suất chuyển $k$ bước:}
\[
  [\PP(Z_{n+k} = j \mid Z_n = i)]_{i,j\in S} = [P^k_{i,j}]_{i,j\in S} = P^k.
\]
Nghĩa là: xác suất chuyển $k$ bước chính là phần tử tương ứng của lũy thừa bậc $k$ của ma trận chuyển.
\end{congthuc}

Quan hệ $P^{m+n} = P^m P^n$ chính là \term{phương trình Chapman--Kolmogorov} (\en{Chapman--Kolmogorov equation}):
\[
  P^{m+n}_{i,j} = \sum_{l \in S} P^m_{i,l}\,P^n_{l,j}, \qquad i,j \in S, \quad m, n \in \N.
\]

\begin{example}[Tính $P^2$ cho chuỗi 3 trạng thái]
\label{ex:P2-three-state}
Xét ma trận chuyển trên $S = \{0, 1, 2\}$:
\[
  P = \begin{bmatrix}
    0 & 1/2 & 1/2 \\
    1/3 & 1/3 & 1/3 \\
    1/2 & 0 & 1/2
  \end{bmatrix}.
\]

Tính $P^2 = P \cdot P$:
\begin{align*}
  P^2_{0,0} &= 0 \cdot 0 + \frac{1}{2} \cdot \frac{1}{3} + \frac{1}{2} \cdot \frac{1}{2} = \frac{1}{6} + \frac{1}{4} = \frac{5}{12}, \\
  P^2_{0,1} &= 0 \cdot \frac{1}{2} + \frac{1}{2} \cdot \frac{1}{3} + \frac{1}{2} \cdot 0 = \frac{1}{6}, \\
  P^2_{0,2} &= 0 \cdot \frac{1}{2} + \frac{1}{2} \cdot \frac{1}{3} + \frac{1}{2} \cdot \frac{1}{2} = \frac{5}{12}.
\end{align*}

Kiểm tra: $P^2_{0,0} + P^2_{0,1} + P^2_{0,2} = 5/12 + 1/6 + 5/12 = 5/12 + 2/12 + 5/12 = 12/12 = 1$. Đúng!

\textbf{Ý nghĩa:} Xuất phát từ trạng thái 0, sau 2 bước, xác suất ở trạng thái 0 là $5/12 \approx 0.42$.
\end{example}

\begin{example}[Tính $P^n$ cho chuỗi Ehrenfest $N = 4$]
\label{ex:Pn-ehrenfest}
Với ma trận $P$ ở Ví dụ~\ref{ex:ehrenfest-N4}, ta tính $P^2$:
\[
  P^2_{0,0} = 0 \cdot 0 + 1 \cdot \frac{1}{4} + 0 + 0 + 0 = \frac{1}{4}.
\]
Nghĩa là: nếu bên trái có 0 bi, sau 2 bước xác suất trở về 0 bi là $1/4$.

Thật vậy: từ 0 phải đi lên 1 (chắc chắn), rồi từ 1 quay về 0 (xác suất $1/4$).
\end{example}

\begin{tomtat}
Các xác suất chuyển $k$ bước được tính bằng cách lũy thừa ma trận chuyển: $P^k$.
Phương trình Chapman--Kolmogorov cho phép phân rã quá trình chuyển thành hai giai đoạn.

\textbf{Mẹo tính toán:} Đối với ma trận $n \times n$ lớn, dùng chéo hóa $P = M D M^{-1}$
sẽ cho $P^k = M D^k M^{-1}$ --- dễ tính hơn nhiều so với nhân ma trận $k$ lần.
\end{tomtat}

% ============================================================================
\section{Chuỗi Markov hai trạng thái (The Two-State Markov Chain)}
\label{sec:two-state}

\textbf{Câu chuyện:} Mô hình đơn giản nhất: thời tiết ngày mai chỉ có ``Nắng'' (0) hoặc ``Mưa'' (1).
Nếu hôm nay nắng, ngày mai mưa với xác suất $a$; nếu hôm nay mưa, ngày mai nắng với xác suất $b$.
Sau nhiều ngày, tỷ lệ ngày nắng/mưa sẽ ổn định --- đây là phân phối giới hạn.

Việc tính ma trận chuyển bậc $n$, tức $P^n$, nói chung là khó. Tuy nhiên, khi số trạng thái bằng 2, tức $S = \{0, 1\}$, ta có thể tính chính xác.

Ma trận chuyển có dạng:
\begin{equation}\label{eq:two-state}
  P = \begin{bmatrix} 1-a & a \\ b & 1-b \end{bmatrix},
\end{equation}
với $a \in [0,1]$ và $b \in [0,1]$.

\textbf{Ý nghĩa:}
\begin{itemize}[nosep]
  \item $\PP(Z_{n+1} = 1 \mid Z_n = 0) = a$: xác suất chuyển từ $0$ sang $1$.
  \item $\PP(Z_{n+1} = 0 \mid Z_n = 0) = 1-a$: xác suất ở lại trạng thái $0$.
  \item $\PP(Z_{n+1} = 0 \mid Z_n = 1) = b$: xác suất chuyển từ $1$ sang $0$.
  \item $\PP(Z_{n+1} = 1 \mid Z_n = 1) = 1-b$: xác suất ở lại trạng thái $1$.
\end{itemize}

Ta loại trừ trường hợp $a = b = 0$ (ma trận đơn vị, chuỗi hằng).

\subsection{Chéo hóa ma trận (Matrix Diagonalization)}

Ma trận $P$ có hai \term{trị riêng} (\en{eigenvalues}):
\[
  \lambda_1 = 1 \qquad \text{và} \qquad \lambda_2 = 1 - a - b,
\]
với các \term{vectơ riêng} (\en{eigenvectors}) tương ứng:
\[
  \begin{bmatrix} 1 \\ 1 \end{bmatrix}
  \qquad \text{và} \qquad
  \begin{bmatrix} -a \\ b \end{bmatrix}.
\]

Do đó $P$ có thể \term{chéo hóa} (\en{diagonalize}): $P = M D M^{-1}$, và:
\begin{equation}\label{eq:Pn-two-state}
  P^n = M D^n M^{-1}
  = \frac{1}{a+b}
  \begin{bmatrix}
    b + a\lambda_2^n & a(1 - \lambda_2^n) \\
    b(1-\lambda_2^n) & a + b\lambda_2^n
  \end{bmatrix}.
\end{equation}

\textbf{Chi tiết tính toán:}
\[
  M = \begin{bmatrix} 1 & -a \\ 1 & b \end{bmatrix}, \quad
  D = \begin{bmatrix} 1 & 0 \\ 0 & 1-a-b \end{bmatrix}, \quad
  M^{-1} = \frac{1}{a+b}\begin{bmatrix} b & a \\ -1 & 1 \end{bmatrix}.
\]

\subsection{Phân phối giới hạn (Limiting Distribution)}

Cho $n \to \infty$ trong \eqref{eq:Pn-two-state}, vì $|\lambda_2| = |1-a-b| < 1$ khi $(a,b) \neq (1,1)$:

\begin{dinhly}[(Phân phối giới hạn của chuỗi hai trạng thái)]
\[
  \lim_{n\to\infty} P^n = \frac{1}{a+b}
  \begin{bmatrix} b & a \\ b & a \end{bmatrix},
\]
do đó:
\begin{align}
  \lim_{n\to\infty} \PP(Z_n = 1 \mid Z_0 = 0) &= \lim_{n\to\infty} \PP(Z_n = 1 \mid Z_0 = 1) = \frac{a}{a+b}, \label{eq:limit-1}\\
  \lim_{n\to\infty} \PP(Z_n = 0 \mid Z_0 = 0) &= \lim_{n\to\infty} \PP(Z_n = 0 \mid Z_0 = 1) = \frac{b}{a+b}. \label{eq:limit-0}
\end{align}
\end{dinhly}

\term{Phân phối giới hạn} (\en{limiting distribution}):
\begin{equation}\label{eq:pi-two-state}
  \pi = [\pi_0, \pi_1] = \left[\frac{b}{a+b},\; \frac{a}{a+b}\right].
\end{equation}

\begin{tomtat}
Phân phối giới hạn $\pi$ \textbf{không phụ thuộc vào trạng thái ban đầu} $Z_0$.
Sau thời gian đủ lớn, xác suất ở trạng thái $1$ xấp xỉ $a/(a+b)$,
và ở trạng thái $0$ xấp xỉ $b/(a+b)$.

Phân phối $\pi$ là \term{bất biến} (\en{invariant}) theo $P$, tức $\pi P = \pi$.

\textbf{Ví dụ thời tiết:} Nếu $a = 0.3$ (nắng$\to$mưa) và $b = 0.4$ (mưa$\to$nắng),
thì dài hạn: $\pi_{\text{nắng}} = 0.4/0.7 \approx 57\%$, $\pi_{\text{mưa}} = 0.3/0.7 \approx 43\%$.
\end{tomtat}

\textbf{Trường hợp đặc biệt:} Khi $a + b = 1$ (hay $a = 1-b$):
\[
  P^n = \begin{bmatrix} b & a \\ b & a \end{bmatrix} = P, \qquad n \in \N,
\]
và $\PP(Z_n = 1 \mid Z_0 = 0) = \PP(Z_n = 1 \mid Z_0 = 1) = a$ với mọi $n \geq 1$.

\textbf{Trường hợp $a = b = 1$:} Giới hạn $\lim_{n\to\infty} P^n$ không tồn tại vì:
\[
  P^n = \begin{cases}
    \begin{bmatrix} 1 & 0 \\ 0 & 1 \end{bmatrix}, & n = 2k, \\[6pt]
    \begin{bmatrix} 0 & 1 \\ 1 & 0 \end{bmatrix}, & n = 2k+1.
  \end{cases}
\]
Chuỗi ``dao động'' mãi mãi --- đây là chuỗi \term{tuần hoàn} với chu kỳ $d = 2$.

\subsection{Tốc độ hội tụ (Rate of Convergence)}

Từ \eqref{eq:Pn-two-state}, sai số giữa $P^n_{i,j}$ và $\pi_j$ là:
\[
  |P^n_{i,j} - \pi_j| \leq C \cdot |\lambda_2|^n = C \cdot |1 - a - b|^n.
\]

Sai số giảm theo cấp số nhân. Hội tụ nhanh khi $|1-a-b|$ nhỏ, tức $a + b$ gần $1$.

\begin{example}[So sánh tốc độ hội tụ]
\begin{itemize}[nosep]
  \item $a = 0.5, b = 0.5$: $|\lambda_2| = |1-1| = 0$. Hội tụ ngay sau 1 bước!
  \item $a = 0.1, b = 0.1$: $|\lambda_2| = 0.8$. Hội tụ chậm (cần $\sim 30$ bước).
  \item $a = 0.01, b = 0.01$: $|\lambda_2| = 0.98$. Rất chậm (cần $\sim 300$ bước).
\end{itemize}
Trực quan: khi $a, b$ nhỏ, chuỗi ``dính'' ở mỗi trạng thái lâu, nên mất nhiều thời gian mới ``quên'' trạng thái ban đầu.
\end{example}

% ============================================================================
\section{Bài tập (Exercises)}
\label{sec:ch4-exercises}

\begin{exercise}
Trong mô hình đơn giản hóa của một chương trình truyền hình, người chơi đã thắng $k$ đô la,
sẽ có $k+1$ đô la ở lượt tiếp theo với xác suất $q$, hoặc bị loại (về $0$) với xác suất $p = 1-q$.
Giả sử người chơi bắt đầu với $1$ đô la. Hãy mô hình hóa tiền thắng sau $n$ lượt chơi
bằng chuỗi Markov và viết ma trận chuyển tương ứng.
\end{exercise}

\begin{exercise}
Xét chuỗi Markov trên $S = \{0, 1, 2\}$ với ma trận chuyển:
\[
  P = \begin{bmatrix}
    1/2 & 1/4 & 1/4 \\
    1/3 & 1/3 & 1/3 \\
    1/4 & 1/4 & 1/2
  \end{bmatrix}.
\]
\begin{enumerate}[label=(\alph*)]
  \item Tính $P^2$ (nhân ma trận).
  \item Tính $\PP(Z_2 = 0 \mid Z_0 = 1)$.
  \item Tính $\PP(Z_3 = 2 \mid Z_0 = 0)$ (gợi ý: dùng $P^3 = P^2 \cdot P$).
\end{enumerate}
\end{exercise}

\begin{exercise}
Cho chuỗi Ehrenfest với $N = 4$ viên bi (Ví dụ~\ref{ex:ehrenfest-N4}).
\begin{enumerate}[label=(\alph*)]
  \item Bắt đầu với 4 bi bên trái. Tính xác suất có 2 bi bên trái sau 2 bước.
  \item Chuỗi này có trạng thái hấp thụ không? Giải thích.
  \item Tìm phân phối dừng $\pi$ (gợi ý: so sánh với phân phối nhị thức $B(4, 1/2)$).
\end{enumerate}
\end{exercise}

\begin{exercise}[Chapman--Kolmogorov trực tiếp]
Xét chuỗi Markov hai trạng thái với $a = 0.4$, $b = 0.6$.
\begin{enumerate}[label=(\alph*)]
  \item Tính $P^2$ trực tiếp bằng nhân ma trận.
  \item Tính $P^2$ bằng công thức \eqref{eq:Pn-two-state} với $n = 2$. Xác minh kết quả trùng nhau.
  \item Tìm $n$ nhỏ nhất sao cho $|P^n_{0,1} - \pi_1| < 0.01$.
\end{enumerate}
\end{exercise}

\begin{exercise}
Chứng minh rằng nếu $(S_n)_{n\geq 0}$ là bước đi ngẫu nhiên với $S_n = X_1 + \cdots + X_n$
và $(X_k)$ độc lập cùng phân phối, thì $(S_n)$ thỏa mãn tính chất Markov.
(Gợi ý: sử dụng tính độc lập của $X_{n+1}$ với $(X_1, \ldots, X_n)$.)
\end{exercise}
